% Français 9115, batch 7 — Gallica views 121–140.
% Pass: Opus 5 (claude-opus-5), 2026-09-15 — first pass, unchecked against
% the views by a human.
%
% What the twenty views hold. Two things, and neither is the memoir on mean
% curvature that batch 5 left in drafts. Views 121 to 134 are printed pages
% of a memoir on the equation $4X = Y^2 \pm nZ^2$ — running heads « MÉMOIRE »
% on the left and « SUR UNE QUESTION D'ANALYSE. » on the right — each page
% printed on one side only and pasted onto a guard leaf, so that every page
% of text shows through, reversed, on the view beside it. No author or title
% is printed on these pages; batch 6 (transcribed in parallel, read after
% this pass had read its views) found the title page at view 115: Legendre's
% « Mémoire sur la détermination des fonctions Y et Z… », read at the
% Académie on 11 October 1830. The piece is already running when the batch
% opens: view 120 is p. 86, and view 121 (p. 87) finishes the table of
% coefficients of Exemple I mid-column. After p. 88 the printed pagination
% starts again at 9 while the text reads on without a gap; the signatures
% (« 2 » on p. 9, « 3 » on p. 17, against « 11 » on p. 81 in batch 6) make
% pp. 9–19 look like gatherings of a separately paginated printing of the
% same memoir, pp. 87–88 the end of the volume's gathering 11 — an inference
% from the marks, not something the leaves say. Then views 136 to 140 are
% three loose leaves in manuscript, working notes on the catenary and the
% elastic surface, headed « Idées à vérifier »; the hand was not compared
% stroke by stroke with the drafts of batches 5 and 6, and is not ascribed
% here. Eighteen views carry a \page{}; view
% 122 is a second exposure of the page already on view 121 and gets only a
% note; 135 and 139, the blank backs of written leaves, are not transcribed.
%
% The printed memoir.
%
%   v. 121     p. 87. End of Example I ($n = 31$): the last coefficients
%              $A_9 \ldots A_{15}$, the remark that they repeat because the
%              equation must not change when $1/x$ is put for $x$, and the
%              polynomials $Y$ and $Z$. Example II, $n = 37$, opens with the
%              table of Legendre symbols $(k/n)$ and stops at « Delà ».
%   v. 122     The same p. 87 photographed a second time. Not transcribed.
%   v. 123     p. 88. Example II carried out: the sums $S$, the coefficients
%              $A_1 \ldots A_{10}$ stopping at the middle term, $Y$ and $Z$
%              for $4X = Y^2 - 37Z^2$; from $n = 41$, and more so at $n = 61$,
%              coefficients exceed $\frac{1}{2}n$. « Voici les résultats de
%              ces calculs qui feront suite au tableau de l'article 512. »
%   v. 124     p. 9 — the pagination starts again, but the text reads on
%              without a gap: the table « Valeurs des polynomes Y et Z » for
%              $n = 31$, 37, 41, 61; then $n = 2521$ as an example of a large
%              prime.
%   v. 125     p. 10. For $n = 2521$ the first coefficients ($A_7 =
%              2633212 + 69292\sqrt{n}$), with $m = 1260$; the growth stops
%              at the middle term because coefficients equidistant from the
%              ends must be equal; with the same sign or with opposite signs?
%   v. 126–128 pp. 11–13. The signs settled: for $n = 4i - 1$, $Y$
%              antisymmetric and $Z$ symmetric (any other choice makes
%              $x = \pm 1$ give $4n = Y^2$ or $4 = nZ^2$); for $n = 4i + 1$
%              both symmetric, the middle coefficients being non-zero.
%   v. 129–131 pp. 14–16. General formulas for $A_1 \ldots A_4$ in the four
%              cases $n = 24\lambda + 1, 5, 13, 17$, carried over to
%              $4i - 1$ by changing the signs of $\lambda$ and of $n$, with
%              the table for $24\lambda - 1, 5, 13, 17$; the first terms of
%              $Y$, $Z$ for $n = 24\lambda - 1$.
%   v. 132–134 pp. 17–19. $1 - 11\lambda + 3\lambda^2$ exceeds $n$ from
%              $n = 311$; the subdivision into sixteen forms $120\mu + r$;
%              for $n = 120\mu + 61$, $A_5$ and $A_6$, and $75\mu^3 > n^2$
%              once $n > 23040$; a further subdivision $840\nu + \ldots$ is
%              declined; $Y$ and $Z$ to $x^{m-6}$ for $120\mu \pm 61$. A
%              printed rule closes p. 19 and, apparently, the piece.
%
% The notes in manuscript.
%
%   v. 136–137 One leaf, both sides, headed « Idées à vérifier ». Whether,
%              in a heavy surface of equal thickness, each normal section
%              is a catenary. With $R$ the radius of the sphere of mean
%              curvature the equation of equilibrium is written
%              $ge + T.\frac{2}{R} = 0$; the section at $45^\circ$ would keep
%              its form if isolated only were the tension halved. A first
%              argument by the supports $B'$, $B''$, $B'''$ of a figure, and
%              an objection the writer says could not be removed
%              (« je n'ai pas trouvé la manière de détruire cette
%              objection »); overleaf a second way — a surface on a circular
%              frame, the tension split between two normal planes, whence
%              $g''e + T'\frac{1}{R} = 0$ and a catenary — and the case of
%              unequal supports, left unfinished.
%   v. 138     The writer has looked for what concerns the catenary « à cause de son
%              rapport avec la voûte » and found it only in the « histoire des
%              mathématiques », the « encyclopédie » and « le mémoire de M
%              poisson sur la surface élastique »; quotes d'Alembert's
%              article (a curve pressed by a normal force is in equilibrium
%              if the force is inversely as the radius of curvature) and
%              writes out the surface version: the half of the force inversely
%              as the radius of the sphere of mean curvature.
%   v. 140     « Voila ce qui doit être examiné »: the element
%              $dxdy\sqrt{1+p^2+q^2}$, its projection $dxdy$, and the
%              variation of $\sqrt{1+p^2+q^2}$ integrated by parts to
%              $\frac{2}{R}\delta z$, where the page stops.
%
% The numbers on the leaves, and why no \folio{} is written. The ink series
% of batches 1 to 5 (40–49 at batch 5) runs on, one number to a guard leaf,
% top right of the side that faces up: 59 at views 121 and 122 (the same
% stroke on both exposures), 60 at 124, 61 at 126, 62 at 128, « 63. » at 130,
% 64 at 132, 65 at 134, 66 at 135 (top left, on the blank side), 67 at 136,
% 68 at 138, 69 at 140. The 67 is in the heavier ink of the text beside it;
% 68 and 69 are in a finer, paler stroke than the text of their leaves, but
% they continue the same series without a break, so they are recorded with
% it. On views 123, 125, 127, 129, 131 and 133 what looks like a number top
% left is the recto's number showing through, reversed — a mirror-image crop
% superposes it on the recto's figure — and on 134 the 66 of view 135 shows
% through above the 65. The printed pages carry their own pagination, 87
% and 88 and then 9 to 19, and printer's signatures at the foot: « 2 » on
% p. 9 (v. 124), « 2. » on p. 11 (v. 126), « 3 » on p. 17 (v. 132). No
% pencilled foliation of the Bibliothèque was found; no \folio{}. The
% Bibliothèque impériale stamp is on view 136.
%
% Notation. The printed pages are transcribed with their own symbols: the
% radical written with a bar-less « V » is set $\sqrt{}$; Legendre's symbol
% is set $\left(\frac{k}{n}\right)$; « etc. » inside formulas as \text{etc.}.
% Braces joining two lines of a polynomial are set with \left\{ and an
% array. In the manuscript the writer's variation sign, a « d » with a
% curled stem, is set $\delta$; the « e » after $g$ in the equilibrium
% equations is kept as the letter written. The writer's underlining is \emph{};
% \uncertain{} renders as an underline too, so the header's word decides.
% Printing slips are left as printed and noted: a « + » where the pattern
% wants « − » in form II of $Z$ (v. 126), $b^{i-1}$ and a stray sign in the
% last value of $Z$ (v. 127), the equation written $Y^2 + nZ^2$ on p. 12 and
% $Y^2 - nZ^2$ on p. 13 for the same case $n = 4i + 1$ (v. 127–128).
%
% What could not be read. Nothing on the printed pages. In the notes, no
% word is left as \ill{}; six readings are offered with doubt and marked:
% « conserve » (v. 136), « Cela » (v. 136), « complémens » (v. 136),
% « élévation » (v. 137, written rather « ébration »), and the struck « les »
% (v. 137) and « je crois » (v. 140), both under their deletions. The primes
% on the figure's letters at view 136 are small and crowded and are read as
% the figure requires. Nothing was guessed to fill a gap.
%
% Nothing here was checked against any published transcription or edition.
\documentclass[11pt,a4paper]{article}
% The preamble shared by every transcription of the mathematicians' manuscripts in Gallica.
%
% It rests on one idea: **uncertainty compounds, it does not resolve.** A
% transcription that smooths over a doubtful word has destroyed the only thing
% it was made for — a reader must be able to tell, at every point, what was on
% the page and what was supplied. The apparatus macros below are therefore the
% critical apparatus, not decoration.
%
% The same commands are recognised by `scripts/render.mjs`, which produces the
% site's reading view, and by `scripts/tei.mjs`. A macro added here must be
% added there too, or rendering fails loudly — which is the wanted behaviour.
%
% Comments here are English, like the rest of the repository. The documents
% this preamble sets are French, and that is deliberate: see the skills.

\ProvidesPackage{germain}[2026/09/15 Transcriptions of mathematicians' manuscripts in Gallica]

\usepackage{amsmath,amssymb,amsthm}
\usepackage{mathrsfs}
% Kept from the parent preamble so the renderer's diagram support stays
% available; these manuscripts draw no commutative diagrams, and nothing here uses it.
\usepackage{tikz-cd}
\usepackage[english,french]{babel}
\usepackage{fontspec}

% --- Greek ------------------------------------------------------------------
% The text face stays fontspec's default, Latin Modern, which has no Greek:
% XeTeX drops every glyph it lacks with a line in the log and nothing on the
% page, and the Greek words the manuscripts quote vanished from the PDFs. So
% the Greek and Greek Extended blocks get a XeTeX character class of their own,
% and entering or leaving a run of it switches the family to CMU Serif — the
% same Computer Modern design, with polytonic Greek — and back. Only the
% family changes: a Greek word in \textit or \textbf keeps its shape and series.
% The transitions to and from every other class carry the tokens that class
% has towards an ordinary letter, so French spacing before « ; » or inside
% « » still holds after a Greek word. No grouping, so no brace or \endgroup
% can come out unbalanced; maths is left alone, since XeTeX inserts nothing
% there. Kept identical in `exercices.sty`.
\makeatletter
\newfontfamily\germainGreekfont{cmun}[
  Extension=.otf, NFSSFamily=germaingreek,
  UprightFont=*rm, ItalicFont=*ti, SlantedFont=*sl,
  BoldFont=*bx, BoldItalicFont=*bi, BoldSlantedFont=*bl]
\newXeTeXintercharclass\germain@greek
\def\germain@greekrange#1#2{%
  \@tempcnta=#1\relax
  \loop
    \XeTeXcharclass\@tempcnta=\germain@greek
    \ifnum\@tempcnta<#2\advance\@tempcnta\@ne\repeat}
\germain@greekrange{"0370}{"03FF}
\germain@greekrange{"1F00}{"1FFF}
\def\germain@greekprev{\rmdefault}
\def\germain@greekin{%
  \edef\germain@greekprev{\f@family}\fontfamily{germaingreek}\selectfont}
\def\germain@greekout{\fontfamily{\germain@greekprev}\selectfont}
\def\germain@greekjoin{%
  \XeTeXinterchartokenstate=\@ne
  \@tempcnta=\z@
  \loop
    \ifnum\@tempcnta=\germain@greek\else
      \XeTeXinterchartoks\germain@greek\@tempcnta=\expandafter{%
        \expandafter\germain@greekout\the\XeTeXinterchartoks\z@\@tempcnta}%
      \XeTeXinterchartoks\@tempcnta\germain@greek=\expandafter{%
        \the\XeTeXinterchartoks\@tempcnta\z@\germain@greekin}%
    \fi
    \ifnum\@tempcnta<4095\advance\@tempcnta\@ne\repeat}
% babel-french sets its own classes, and saves and restores the interchar
% state, each time a language is selected: join again after it does.
\addto\extrasfrench{\germain@greekjoin}
\addto\extrasenglish{\germain@greekjoin}
\AtBeginDocument{\germain@greekjoin}
\makeatother

\usepackage{geometry}
\geometry{a4paper,margin=25mm,marginparwidth=28mm}
\usepackage{marginnote}
\usepackage{xcolor}
\usepackage[normalem]{ulem}
\setlength{\emergencystretch}{3em}
\usepackage{hyperref}
\hypersetup{colorlinks=true,linkcolor=black,urlcolor=black,pdfborder={0 0 0}}

% --- Batch metadata ---------------------------------------------------------
% Taken as they stand by the rendering script, which shows them at the head of
% the reading view. They fix what the file claims to cover. The macro names are
% the parent project's, so that the scripts stay shared; \folder here means the
% volume's slug — `fr-9115` — and \pages counts Gallica views.

\newcommand{\folder}[1]{\gdef\grothFolder{#1}}
\newcommand{\batch}[1]{\gdef\grothBatch{#1}}
\newcommand{\pages}[2]{\gdef\grothFirst{#1}\gdef\grothLast{#2}}
\newcommand{\foldertitle}[1]{\gdef\grothFoldertitle{#1}}

% The holder's own shelfmark, printed as they write it: « Français 9115 ».
\newcommand{\shelfmark}[1]{\gdef\grothShelfmark{#1}}

% The dating, copied verbatim from the holder's catalogue — never inferred
% here. « XIXe siècle » is the BnF's; a pass that dates a leaf from its content
% says so in a \note{}, as its own inference.
\newcommand{\dating}[1]{\gdef\grothDating{#1}}

% The status notice, printed in the title block and repeated as a diagonal
% watermark on every page of the PDF. The manuscripts are in the public
% domain; what this line declares is not a right but a quality — an
% unverified first pass by a machine. The skills write the line; a file
% without it gets no watermark, so leaving it out is visible in the output.
\newcommand{\watermark}[1]{\gdef\grothWatermark{#1}}

\gdef\grothFolder{?}\gdef\grothBatch{}\gdef\grothFirst{?}\gdef\grothLast{?}%
\gdef\grothFoldertitle{}\gdef\grothShelfmark{}\gdef\grothDating{}\gdef\grothWatermark{}

% --- The résumé -------------------------------------------------------------
\newenvironment{resume}
  {\small\setlength{\parskip}{0.45em}}
  {\par\medskip\hrule\medskip}

% --- Keywords ---------------------------------------------------------------
% In English, closing the modernised reading's résumé: the modern vocabulary
% under which to search for what the leaves construct. The manifest extracts
% them to tag the volume — this line is the single source of the tags.
\newcommand{\keywords}[1]{\par\smallskip\noindent
  {\footnotesize\textsc{Keywords} --- \textit{#1}}\par}

% --- The critical apparatus -------------------------------------------------

% The Gallica view number, in the margin. This is the anchor that lets a
% disagreement over a formula be settled by looking at *one* image rather than
% a volume of 750. The site uses it to turn the facsimile as the transcription
% is scrolled. A view is one scanned image: usually one side of a leaf.
\newcommand{\page}[1]{%
  \marginnote{\footnotesize\color{gray!70}\textsf{v.\,#1}}%
  \label{page:#1}\ignorespaces}

% The BnF's pencilled foliation, where the leaf carries it and a pass can read
% it: « 198r ». This is what the literature cites — « ff. 198r–208v » — and
% the one line that turns a citation into a view. Recorded, never inferred:
% a folio a pass cannot read is not written.
\newcommand{\folio}[1]{%
  \marginnote{\footnotesize\color{gray!50}\textsf{f.\,#1}}\ignorespaces}

% The modernised reading's coarser counterpart to \page{}: which views a
% section is drawn from.
\newcommand{\pagerange}[2]{%
  \marginnote{\footnotesize\color{gray!70}\textsf{v.\,#1--#2}}%
  \ignorespaces}

% Illegible: never guessed.
\newcommand{\ill}{\textcolor{red!60!black}{[\,\ldots\,]}}

% A doubtful reading: the word is given, and flagged as uncertain.
\newcommand{\uncertain}[1]{\uline{#1}}

% An editorial addition — a missing letter, an implied word.
\newcommand{\add}[1]{\textcolor{blue!50!black}{[#1]}}

% Struck out by the writer. What was crossed out is part of the document.
\newcommand{\struck}[1]{\sout{#1}}

% The transcriber's note, on the state of the page or the mathematics.
\newcommand{\note}[1]{\par\smallskip{\small\color{gray!60!black}\textit{[#1]}}\par\smallskip}

% A marginal note of hers, distinct from the body of the text.
\newcommand{\marginal}[1]{\par\smallskip\noindent\hspace{1em}%
  {\small\color{gray!60!black}#1}\par\smallskip}

% --- Title ------------------------------------------------------------------
% The title block is French, like the documents themselves.

\AddToHook{shipout/background}{%
  \ifx\grothWatermark\empty\else
    \begin{tikzpicture}[remember picture, overlay]
      \node[rotate=52, scale=3.4, align=center, text=gray!18,
            font=\sffamily\bfseries]
        at (current page.center) {\grothWatermark};
    \end{tikzpicture}%
  \fi}

\AtBeginDocument{%
  \begin{center}
    {\large\bfseries Manuscrits de mathématiciens (Gallica) ---
      \ifx\grothShelfmark\empty\grothFolder\else\grothShelfmark\fi}\\[2pt]
    {\itshape\grothFoldertitle}\\[4pt]
    {\small vues \grothFirst--\grothLast
      \ifx\grothBatch\empty\else\ (lot \grothBatch)\fi}\\[2pt]
    \ifx\grothDating\empty\else
      {\small\color{gray!60!black}Datation du catalogue : \grothDating}\\[2pt]
    \fi
    \ifx\grothWatermark\empty\else
      {\small\bfseries\color{red!45!gray}\grothWatermark}\\[2pt]
    \fi
    {\footnotesize\color{gray!60!black}Transcription établie d'après les images IIIF de Gallica.
     Source gallica.bnf.fr / Bibliothèque nationale de France.\\
     Les lectures incertaines sont soulignées, les illisibles notées [\,\ldots\,],
     les ajouts entre crochets ; \textsf{v.} numérote les vues de Gallica,
     \textsf{f.} la foliotation au crayon de la BnF.}
  \end{center}
  \vspace{1em}}

\endinput


\folder{fr-9115}
\shelfmark{Français 9115}
\batch{7}
\pages{121}{140}
\dating{1801-1900}
\watermark{Lecture automatique\\non vérifiée}
\foldertitle{Papiers de Sophie GERMAIN. — Recueil de dissertations et problèmes mathématiques et physiques.}

\begin{document}

\note{Numérotation des feuillets. La série à l'encre des lots précédents se
poursuit, un nombre par feuillet de montage : 59 (vues 121 et 122, seconde
prise de vue de la même page), 60 (vue 124), 61 (vue 126), 62 (vue 128),
« 63. » (vue 130), 64 (vue 132), 65 (vue 134), 66 (vue 135, au côté blanc),
67 (vue 136), 68 (vue 138), 69 (vue 140). Ce que l'on croit lire en haut à
gauche des vues 123, 125, 127, 129, 131 et 133 est le nombre du recto vu en
transparence. Les pages imprimées portent leur propre pagination, 87, 88
puis 9 à 19, et au pied les signatures « 2 » (p. 9), « 2. » (p. 11) et
« 3 » (p. 17). Aucune foliotation au crayon de la Bibliothèque n'a été lue ;
aucun « f. » n'est porté. Les vues 135 et 139 sont les dos blancs de
feuillets écrits et ne sont pas transcrites ; la vue 122 répète la page de la
vue 121 et n'est pas transcrite une seconde fois.}

\page{121}
\note{page imprimée, titre courant « SUR UNE QUESTION D'ANALYSE. » et
numéro de page 87 ; le nombre 59 à l'encre en haut à droite. Le texte
s'arrête au mot « Delà », au tiers inférieur de la page, dont le reste est
blanc ; le texte qu'on devine en transparence est celui de la page 88,
collée de l'autre côté du même feuillet de montage (vue 123), et n'est
pas transcrit.}

\[
\begin{array}{l}
A_9 = 3 + \sqrt{-n} \\
A_{10} = -8 \\
A_{11} = -2 - 2\sqrt{-n} \\
A_{12} = 11 - \sqrt{-n} \\
A_{13} = 7 + \sqrt{-n} \\
A_{14} = -1 + \sqrt{-n} \\
A_{15} = -2.
\end{array}
\]

C'est une suite nécessaire de ce que l'équation $4X = Y^2 + nZ^2$ ne doit
pas changer de forme en mettant $\frac{1}{x}$ à la place de $x$, propriété
sur laquelle nous reviendrons ci-après.

Connaissant les coefficients A qui servent à composer la fonction
\[
Y + Z\sqrt{-n} = 2x^{15} + A_1 x^{14} + A_2 x^{13} + A_3 x^{12} + \text{etc.},
\]
on en déduira les valeurs des polynomes Y et Z, savoir :
\[
\begin{array}{l}
Y = 2x^{15} + x^{14} - 7x^{13} - 11x^{12} + 2x^{11} + 8x^{10} - 3x^9 - 5x^8 \\
\qquad + 5x^7 + 3x^6 - 8x^5 - 2x^4 + 11x^3 + 7x^2 - x - 2 \\
Z = x^{14} + x^{13} - x^{12} - 2x^{11} + x^9 - x^8 - x^7 + x^6 - 2x^4 - x^3 \\
\qquad + x^2 + x.
\end{array}
\]

\subsection{Exemple II.}

Soit $n = 37$, on aura
\[
\begin{array}{l}
\left(\frac{2}{n}\right) = -1,\ \left(\frac{3}{n}\right) = 1,\ \left(\frac{4}{n}\right) = 1,\ \left(\frac{5}{n}\right) = -1,\ \left(\frac{6}{n}\right) = -1 \\
\left(\frac{7}{n}\right) = 1,\ \left(\frac{8}{n}\right) = -1,\ \left(\frac{9}{n}\right) = 1,\ \left(\frac{10}{n}\right) = 1.
\end{array}
\]

Delà

\page{122}
\note{seconde prise de vue de la même page imprimée 87 que la vue 121 : même
texte, arrêté au même « Delà », et même nombre 59 à l'encre, dont le tracé,
comparé sur les deux vues, est identique trait pour trait. Le feuillet est
seulement photographié sous un autre cadrage. Le texte n'est pas transcrit
une seconde fois.}

\page{123}
\note{page imprimée, titre courant « MÉMOIRE » et numéro de page 88,
imprimée d'un seul côté et collée au verso du feuillet de montage qui porte
au recto la page 87 (vue 121). Le « 59 » qu'on croit lire en haut à gauche
n'est pas écrit de ce côté : c'est le nombre à l'encre du recto vu en
transparence ; l'image retournée en miroir se superpose trait pour trait au
59 de la vue 121. Le texte fantôme qui se lit à l'envers est de même celui
de la page 87.}

\[
\begin{array}{l}
S_1 = S_3 = S_4 = S_7 = S_9 = S_{10} = -\frac{1}{2} - \frac{1}{2}\sqrt{n} \\
S_2 = S_5 = S_6 = S_8 = -\frac{1}{2} + \frac{1}{2}\sqrt{n},
\end{array}
\]
et la substitution de ces valeurs dans les équations (1) donnera
\[
\begin{array}{l}
A_1 = 1 + \sqrt{n},\ A_2 = 10,\ A_3 = -4 + 2\sqrt{n},\ A_4 = 15 - \sqrt{n}, \\
A_5 = -5 + 3\sqrt{n},\ A_6 = 17 - \sqrt{n},\ A_7 = -8 + 2\sqrt{n}, \\
A_8 = 11 - \sqrt{n},\ A_9 = -4 + 2\sqrt{n},\ A_{10} = 11 - \sqrt{n}.
\end{array}
\]

Le terme moyen étant $A_9$, il est inutile d'aller plus loin, et les
coefficients qui suivent $A_9$ seront égaux à ceux qui le précèdent à
intervalles égaux, de sorte qu'on aura
\[
A_{10} = A_8 = 11 - \sqrt{n},\quad A_{11} = A_7 = -8 + 2\sqrt{n},\ \text{etc.};
\]
et on en déduit les valeurs suivantes des fonctions Y et Z,
\[
\begin{array}{l}
Y = 2x^{18} + x^{17} + 10x^{16} - 4x^{15} + 15x^{14} - 5x^{13} + 17x^{12} - 8x^{11} \\
\qquad + 11x^{10} - 4x^9 + 11x^8 - 8x^7 + 17x^6 - 5x^5 + 15x^4 - 4x^3 \\
\qquad + 10x^2 + x + 2 \\
Z = x^{17} + 2x^{15} - x^{14} + 3x^{13} - x^{12} + 2x^{11} - x^{10} + 2x^9 \\
\qquad - x^8 + 2x^7 - x^6 + 3x^5 - x^4 + 2x^3 + x,
\end{array}
\]
lesquelles satisfont à l'équation $4X = Y^2 - 37Z^2$.

Dans ces deux exemples les coefficients des fonctions Y et Z sont encore
plus petits que $\frac{1}{2}n$ ; mais à compter de $n = 41$, on trouve des
coefficients plus grands, ce qui devient encore plus marqué, lorsque
franchissant un plus grand intervalle on prend $n = 61$. Voici les
résultats de ces calculs qui feront suite au tableau de l'article 512.

\page{124}
\note{page imprimée, titre courant « SUR UNE QUESTION D'ANALYSE. » et
numéro de page 9 — et non 89 : à partir d'ici l'imprimé porte une autre
pagination, bien que le texte fasse suite à la page 88 de la vue 123, dont
il donne le tableau annoncé. Le nombre 60 est à l'encre en haut à droite ;
au pied de la page, la signature d'imprimeur « 2 ». Quelques taches
d'encre dans la marge et sur les mots « les » et « premier ».}

\textit{Valeurs des polynomes} Y \textit{et} Z.
\[
\begin{array}{r|l}
n & \\ \hline
31 & Y = 2x^{15} + x^{14} - 7x^{13} - 11x^{12} + 2x^{11} + 8x^{10} - 3x^9 - 5x^8 \\
   & \qquad + 5x^7 + 3x^6 - 8x^5 - 2x^4 + 11x^3 + 7x^2 - x - 2 \\
   & Z = x^{14} + x^{13} - x^{12} - 2x^{11} + x^9 - x^8 - x^7 + x^6 - 2x^4 - x^3 + x^2 + x \\ \hline
37 & Y = 2x^{18} + x^{17} + 10x^{16} - 4x^{15} + 15x^{14} - 5x^{13} + 17x^{12} - 8x^{11} + 11x^{10} - 4x^9 \\
   & \qquad + 11x^8 - 8x^7 + 17x^6 - 5x^5 + 15x^4 - 4x^3 + 10x^2 + x + 2 \\
   & Z = x^{17} + 2x^{15} - x^{14} + 3x^{13} - x^{12} + 2x^{11} - x^{10} + 2x^9 \\
   & \qquad - x^8 + 2x^7 - x^6 + 3x^5 - x^4 + 2x^3 + x \\ \hline
41 & Y = 2x^{20} + x^{19} + 11x^{18} + 16x^{17} + 14x^{16} + 29x^{15} + 30x^{14} \\
   & \qquad + 22x^{13} + 36x^{12} + 34x^{11} + 20x^{10} + 34x^9 + 36x^8 + 22x^7 \\
   & \qquad + 30x^6 + 29x^5 + 14x^4 + 16x^3 + 11x^2 + x + 2 \\
   & Z = x^{19} + x^{18} + 2x^{17} + 4x^{16} + 3x^{15} + 4x^{14} + 6x^{13} + 4x^{12} + 4x^{11} \\
   & \qquad + 6x^{10} + 4x^9 + 4x^8 + 6x^7 + 4x^6 + 3x^5 + 4x^4 + 2x^3 + x^2 + x \\ \hline
61 & Y = 2x^{30} + x^{29} + 16x^{28} - 7x^{27} + 32x^{26} - 20x^{25} + 63x^{24} - 33x^{23} \\
   & \qquad + 72x^{22} - 54x^{21} + 89x^{20} - 62x^{19} + 88x^{18} - 89x^{17} + 95x^{16} \\
   & \qquad - 81x^{15} + \text{etc.} \\
   & Z = x^{29} + 3x^{27} - 2x^{26} + 6x^{25} - 3x^{24} + 9x^{23} - 6x^{22} + 10x^{21} \\
   & \qquad - 7x^{20} + 12x^{19} - 10x^{18} + 11x^{17} - 11x^{16} + 13x^{15} - \text{etc.}
\end{array}
\]

On voit dans ce tableau que pour le nombre 61 plusieurs des coefficients
de Y sont plus grands que $n$, mais pour avoir une idée de la progression
suivant laquelle les coefficients augmentent lorsque $n$ est un nombre
premier un peu grand, prenons par exemple $n = 2521$, nous aurons
\[
\left(\frac{2}{n}\right) = 1,\ \left(\frac{3}{n}\right) = 1,\ \left(\frac{4}{n}\right) = 1,\ \left(\frac{5}{n}\right) = 1,\ \left(\frac{6}{n}\right) = 1,\ \left(\frac{7}{n}\right) = 1,\ \text{etc.},
\]
ce qui donne

\page{125}
\note{page imprimée, titre courant « MÉMOIRE » et numéro de page 10, collée
au verso du feuillet de montage dont le recto porte la page 9 (vue 124). Le
« 60 » à rebours en haut à gauche est le nombre du recto vu en transparence.
Taches d'encre dans la marge de gauche.}

\[
S_1 = S_2 = S_3 = S_4 = S_5 = S_6 = S_7 = -\frac{1}{2} - \frac{1}{2}\sqrt{n}.
\]
Substituant ces valeurs dans les équations (1), on en tire
\[
\begin{array}{ll}
A_1 = 1 + \sqrt{n} & A_5 = 84001 + 3613\sqrt{n} \\
A_2 = 631 + \sqrt{n} & A_6 = 836032 + 10312\sqrt{n} \\
A_3 = 946 + 106\sqrt{n} & A_7 = 2633212 + 69292\sqrt{n} \\
A_4 = 34231 + 211\sqrt{n} & \text{etc.}
\end{array}
\]

Ainsi les premiers termes des fonctions Y et Z seront, en faisant
$1260 = m$
\[
\begin{array}{l}
Y = 2x^m + x^{m-1} + 631x^{m-2} + 946x^{m-3} \\
\qquad + 34231x^{m-4} + 84001x^{m-5} \\
\qquad + 836032x^{m-6} + 2633212x^{m-7} + \text{etc.} \\
Z = x^{m-1} + x^{m-2} + 106x^{m-3} + 211x^{m-4} \\
\qquad + 3613x^{m-5} + 10312x^{m-6} \\
\qquad + 69292x^{m-7} + \text{etc.}
\end{array}
\]

On voit avec quelle rapidité les coefficients croissent dans l'une et
l'autre fonction, mais cet accroissement n'a lieu que jusqu'au terme
moyen ; ensuite il fait place à un décroissement semblable, puisque les
coefficients doivent être égaux à égale distance des extrêmes.

C'est en effet la loi à laquelle sont assujéties les fonctions Y et Z,
puisque la substitution de $x^{-1}$ à la place de $x$ dans l'équation
$4X = Y^2 - nZ^2$, n'y doit apporter aucun changement, après qu'on a fait
disparaître le dénominateur commun. On voit par là que dans chaque fonction
les coefficients des termes extrêmes ou également éloignés des extrêmes,
doivent être égaux ; mais ces coefficients sont-ils précédés du même signe
ou de signes différents ? c'est ce qu'il faut examiner.

\page{126}
\note{page imprimée, titre courant « SUR UNE QUESTION D'ANALYSE. » et
numéro de page 11 ; le nombre 61 à l'encre en haut à droite ; au pied de la
page, la marque « 2. ».}

Soit 1\textsuperscript{o} $n = 4i - 1$ ou $m = \frac{1}{2}(n-1) = 2i - 1$ :
puisque la fonction Z doit rester la même, ou du moins ne faire que
changer de signe, en mettant $x^{-1}$ à la place de $x$ et multipliant le
tout par $x^m$, cette fonction ne pourra avoir que l'une des deux valeurs
suivantes :
\[
\begin{array}{ll}
\text{I.} & Z = \left\{ \begin{array}{l}
x^{m-1} + b_2 x^{m-2} + b_3 x^{m-3} \ldots + b_{i-1} x^i \\
+ x + b_2 x^2 + b_3 x^3 \ldots + b_{i-1} x^{i-1}
\end{array} \right. \\[1ex]
\text{II.} & Z = \left\{ \begin{array}{l}
x^{m-1} + b_2 x^{m-2} + b_3 x^{m-3} \ldots + b_{i-1} x^i \\
- x - b_2 x^2 - b_3 x^3 \ldots + b_{i-1} x^{i-1}.
\end{array} \right.
\end{array}
\]
\note{dans la seconde ligne de la forme II, le dernier terme est bien
précédé du signe $+$ sur la page, là où la suite des signes appelleroit
$-$ ; la ligne correspondante de la forme II de Y, plus bas, porte
$-a_{i-1}x^{i-1}$.}

Or si la seconde valeur avait lieu, la supposition $x = 1$ donnerait
$Z = 0$, et comme en même temps on aurait $X = \frac{x^n - 1}{x - 1} = n$,
l'équation $4X = Y^2 + nZ^2$ deviendrait $4n = Y^2$, équation impossible,
puisque $4n$ n'est point un carré, donc la forme I a lieu nécessairement.

Cela posé cette forme I de la fonction Z doit être combinée avec l'une des
deux formes que pourrait avoir la fonction Y, lesquelles sont
\[
\begin{array}{ll}
\text{I.} & Y = \left\{ \begin{array}{l}
2x^m + x^{m-1} + a_2 x^{m-2} + a_3 x^{m-3} \ldots + a_{i-1} x^i \\
+ 2 + x + a_2 x^2 + a_3 x^3 \ldots + a_{i-1} x^{i-1}
\end{array} \right. \\[1ex]
\text{II.} & Y = \left\{ \begin{array}{l}
2x^m + x^{m-1} + a_2 x^{m-2} + a_3 x^{m-3} \ldots + a_{i-1} x^i \\
- 2 - x - a_2 x^2 - a_3 x^3 \ldots - a_{i-1} x^{i-1}.
\end{array} \right.
\end{array}
\]

Or si la forme I avait lieu, la supposition $x = -1$ donnerait $Y = 0$, et
comme alors on aurait $X = 1$, l'équation $4X = Y^2 + nZ^2$ deviendrait
$4 = nZ^2$, équation impossible ; donc la forme II de Y est celle qui a
lieu nécessairement.

Donc lorsqu'on a $n = 4i - 1$, les fonctions Y et Z qui satisfont à
l'équation $4X = Y^2 + nZ^2$ ont nécessairement la

\page{127}
\note{page imprimée, titre courant « MÉMOIRE » et numéro de page 12, collée
au verso du feuillet dont le recto porte la page 11 (vue 126). Le « 61 » à
rebours en haut à gauche est le nombre du recto vu en transparence.}

forme ici déterminée
\[
\begin{array}{l}
Y = \left\{ \begin{array}{l}
2x^m + x^{m-1} + a_2 x^{m-2} + a_3 x^{m-3} \ldots + a_{i-1} x^i \\
- 2 - x - a_2 x^2 - a_3 x^3 \ldots - a_{i-1} x^{i-1}
\end{array} \right. \\[1ex]
Z = \left\{ \begin{array}{l}
x^{m-1} + b_2 x^{m-2} + b_3 x^{m-3} \ldots + b_{i-1} x^i \\
+ x + b_2 x^2 + b_3 x^3 \ldots + b_{i-1} x^{i-1}.
\end{array} \right.
\end{array}
\]

Le tableau que nous avons donné de ces fonctions pour quelques-unes des
plus simples valeurs de $n$ s'accorde avec ce résultat, mais il était
nécessaire de démontrer l'existence de cette propriété indépendamment de
toute induction.

Soit 2\textsuperscript{o} $n = 4i + 1$ ou $m = 2i$, nous allons faire voir
que la valeur de Z ne peut être que de la forme
\[
Z = \left\{ \begin{array}{l}
x^{m-1} + b_2 x^{m-2} + b_3 x^{m-3} \ldots + b_{i-1} x^{i+1} + b_i x^i \\
+ x + b_2 x^2 + b_3 x^3 \ldots + b_{i-1} x^{i-1}
\end{array} \right.
\]
qui ne changera pas en mettant $x^{-1}$ à la place de $x$ et multipliant
le tout par $x^m$. Car si après le terme moyen $b_i x^i$ tous les signes
changeaient, en sorte qu'on eût
\[
Z = \left\{ \begin{array}{l}
x^{m-1} + b_2 x^{m-2} + b_3 x^{m-3} \ldots + b_{i-1} x^{i+1} + b_i x^i \\
- x - b_2 x^2 - b_3 x^3 \ldots - b_{i-1} x^{i-1},
\end{array} \right.
\]
la substitution de $x^{-1}$ au lieu de $x$, donnerait, après avoir
multiplié le tout par $-x^m$,
\[
Z = \left\{ \begin{array}{l}
x^{m-1} + b_2 x^{m-2} + b_3 x^{m-3} \ldots + b^{i-1} x^{i+1} - b_i x^i \\
- x + b_2 x^2 - b_3 x^3 \ldots - b_{i-1} x^{i-1},
\end{array} \right.
\]
\note{dans cette dernière valeur, l'imprimé porte $b^{i-1}$, l'indice
placé en exposant, et $+b_2x^2$ à la seconde ligne, entre $-x$ et
$-b_3x^3$ : l'un et l'autre sont transcrits tels qu'on les lit sur la
page.}
valeur qui ne pourrait être égale à la précédente qu'en supposant
$b_i = 0$. Mais alors la supposition $x = 1$ donnerait $Z = 0$ et
l'équation $4X = Y^2 + nZ^2$ deviendrait $4n = Y^2$, équation

\page{128}
\note{page imprimée, titre courant « SUR UNE QUESTION D'ANALYSE. » et
numéro de page 13 ; le nombre 62 à l'encre en haut à droite.}

impossible. Donc la forme que nous avons donnée à la valeur de Z est la
seule admissible.

Maintenant cette valeur de Z doit être combinée avec l'une ou l'autre de
ces deux valeurs de Y
\[
\begin{array}{ll}
\text{I.} & Y = \left\{ \begin{array}{l}
2x^m + x^{m-1} + a_2 x^{m-2} + a_3 x^{m-3} \ldots + a_{i-1} x^{i+1} + a_i x^i \\
+ 2 + x + a_2 x^2 + a_3 x^3 \ldots + a_{i-1} x^{i-1}
\end{array} \right. \\[1ex]
\text{II.} & Y = \left\{ \begin{array}{l}
2x^m + x^{m-1} + a_2 x^{m-2} + a_3 x^{m-3} \ldots + a_{i-1} x^{i+1} + a_i x^i \\
- 2 - x - a_2 x^2 - a_3 x^3 \ldots - a_{i-1} x^{i-1}.
\end{array} \right.
\end{array}
\]

Or si la seconde forme avait lieu, il faudrait qu'elle ne changeât pas, ou
qu'elle changeât seulement de signe, en mettant $x^{-1}$ à la place de $x$
et multipliant tout par $x^m$, ce qui exigerait qu'on eût $a_i = 0$. Mais
cette condition ne peut être remplie, puisqu'on sait qu'aucune des
puissances de $x$ ne peut manquer dans le polynome Y ; d'ailleurs si le
coefficient $a_i$ était zéro, la supposition $x = 1$ dans la seconde forme
donnerait $Y = 0$, et alors l'équation $4X = Y^2 - nZ^2$ deviendrait
$4n = -nZ^2$ équation impossible ; donc par cette double raison la première
forme de Y est la seule admissible ; donc dans le cas de $n = 4i + 1$ les
fonctions Y et Z qui satisfont à l'équation $4X = Y^2 - nZ^2$ seront
toujours des formes suivantes :
\[
\begin{array}{l}
Y = \left\{ \begin{array}{l}
2x^m + x^{m-1} + a_2 x^{m-2} + a_3 x^{m-3} \ldots + a_{i-1} x^{i+1} + a_i x^i \\
+ 2 + x + a_2 x^2 + a_3 x^3 \ldots + a_{i-1} x^{i-1}
\end{array} \right. \\[1ex]
Z = \left\{ \begin{array}{l}
x^{m-1} + b_2 x^{m-2} + b_3 x^{m-3} \ldots + b_{i-1} x^{i+1} + b_i x^i \\
+ x + b_2 x^2 + b_3 x^3 \ldots + b_{i-1} x^{i-1},
\end{array} \right.
\end{array}
\]
résultat qui s'accorde avec tous les exemples compris dans notre tableau.
\note{à la page précédente (vue 127), le même cas $n = 4i + 1$ est
rapporté à l'équation $4X = Y^2 + nZ^2$ ; celle-ci écrit $4X = Y^2 - nZ^2$.
Les deux signes sont ceux de l'imprimé.}

La méthode que nous avons donnée pour calculer les va-

\page{129}
\note{page imprimée, titre courant « MÉMOIRE » et numéro de page 14, collée
au verso du feuillet dont le recto porte la page 13 (vue 128). Le « 62 » à
rebours en haut à gauche est le nombre du recto vu en transparence.}

leurs des fonctions Y et Z, ne laisse rien à désirer quand il s'agit de
l'appliquer à une valeur particulière du nombre premier $n$ ; mais pour
parvenir à quelques résultats généraux, il sera bon de considérer la
valeur de $n$ sous une forme indéterminée qui la rende applicable à une
infinité de nombres premiers. Et parce que les formules relatives aux
nombres premiers $4i - 1$, se déduisent aisément des formules relatives
aux nombres premiers $4i + 1$, nous allons nous occuper exclusivement de
ceux-ci.

Observons d'abord que la forme $4i + 1$ se divise en quatre autres,
savoir :
\[
24\lambda + 1,\quad 24\lambda + 5,\quad 24\lambda + 13,\quad 24\lambda + 17,
\]
ce qui offre quatre cas à considérer.

\textit{Premier cas} $n = 24\lambda + 1$.

Alors on aura $\left(\frac{2}{n}\right) = 1$,
$\left(\frac{3}{n}\right) = \left(\frac{n}{3}\right) = 1$,
$\left(\frac{4}{n}\right)$ ou en général $\left(\frac{c^2}{n}\right) = 1$.
Nous nous arrêterons là, parce que pour aller plus loin, il faudrait
connaître le reste que donne $\lambda$ divisé par 5, ce qui exigerait une
nouvelle subdivision. Au moyen de ces valeurs de
$\left(\frac{4}{n}\right)$, on trouvera
\[
S_1 = S_2 = S_3 = -\frac{1}{2} - \frac{1}{2}\sqrt{n};
\]
ensuite par les équations (1), on aura
\[
\begin{array}{l}
A_1 = 1 + \sqrt{n} \\
A_2 = 1 + 6\lambda + \sqrt{n} \\
A_3 = 1 + 9\lambda + (1 + \lambda)\sqrt{n} \\
A_4 = 1 + 11\lambda + 3\lambda^2 + (1 + 2\lambda)\sqrt{n}.
\end{array}
\]

\page{130}
\note{page imprimée, titre courant « SUR UNE QUESTION D'ANALYSE. » et
numéro de page 15 ; le nombre « 63. » à l'encre en haut à droite.}

\textit{Second cas} $n = 24\lambda + 5$.

Alors on aura $\left(\frac{2}{n}\right) = -1$,
$\left(\frac{3}{n}\right) = \left(\frac{5}{3}\right) = -1$,
$\left(\frac{4}{n}\right) = 1$ ; ce qui donne
\[
S_1 = S_4 = -\frac{1}{2} - \frac{1}{2}\sqrt{n},\quad S_2 = S_3 = -\frac{1}{2} + \frac{1}{2}\sqrt{n},
\]
et par l'application des équations (1), on aura
\[
\begin{array}{l}
A_1 = 1 + \sqrt{n} \\
A_2 = 2 + 6\lambda \\
A_3 = -3\lambda + \lambda\sqrt{n} \\
A_4 = -2\lambda + 3\lambda^2 - \lambda\sqrt{n}.
\end{array}
\]

\textit{Troisième cas}, $n = 24\lambda + 13$.

Alors on aura $\left(\frac{2}{n}\right) = -1$, $\left(\frac{3}{n}\right) = 1$,
$\left(\frac{4}{n}\right) = 1$, ensuite
\[
S_1 = S_3 = S_4 = -\frac{1}{2} - \frac{1}{2}\sqrt{n},\quad S_2 = -\frac{1}{2} + \frac{1}{2}\sqrt{n},
\]
et par l'application des équations (1) on trouve
\[
\begin{array}{l}
A_1 = 1 + \sqrt{n} \\
A_2 = 4 + 6\lambda \\
A_3 = -1 - 3\lambda + (1 + \lambda)\sqrt{n} \\
A_4 = 4 + 8\lambda + 3\lambda^2 - \lambda\sqrt{n}.
\end{array}
\]

\textit{Quatrième cas}, $n = 24\lambda + 17$.

Alors on aura $\left(\frac{2}{n}\right) = 1$, $\left(\frac{3}{n}\right) = -1$,
$\left(\frac{4}{n}\right) = 1$ ; de là
\[
S_1 = S_2 = S_4 = -\frac{1}{2} - \frac{1}{2}\sqrt{n},\quad S_3 = -\frac{1}{2} + \frac{1}{2}\sqrt{n},
\]

\page{131}
\note{page imprimée, titre courant « MÉMOIRE » et numéro de page 16, collée
au verso du feuillet dont le recto porte la page 15 (vue 130). Le « 63. »
à rebours en haut à gauche est le nombre du recto vu en transparence.}

et en vertu des équations (1) ;
\[
\begin{array}{l}
A_1 = 1 + \sqrt{n} \\
A_2 = 5 + 6\lambda + \sqrt{n} \\
A_3 = 7 + 9\lambda + (1 + \lambda)\sqrt{n} \\
A_4 = 4 + 7\lambda + 3\lambda^2 + (2 + 2\lambda)\sqrt{n}.
\end{array}
\]

Ces formules pour les nombres premiers $4i + 1$, considérés sous les quatre
formes $24\lambda + 1$, $+5$, $+13$, $+17$, dont ils sont susceptibles,
s'appliquent aux nombres premiers $4i - 1$, considérés également sous leurs
quatre formes $24\lambda - 1$, $-5$, $-13$, $-17$. Il suffit pour cela de
changer le signe de $\lambda$ dans les formules précédentes et en même temps
le signe de $n$ dans l'équation $4X = Y^2 - nZ^2$ relative aux nombres
premiers $4i + 1$ ; on aura ainsi les valeurs de A qui satisfont à
l'équation $4X = Y^2 + nZ^2$ relative aux nombres premiers $4i - 1$ ; nous
les avons réunies dans le tableau suivant :
\[
\begin{array}{c|c|c|c|c}
n & A_1 & A_2 & A_3 & A_4 \\ \hline
24\lambda - 1 & 1 + \sqrt{-n} & 1 - 6\lambda + \sqrt{-n} & 1 - 9\lambda + (1 - \lambda)\sqrt{-n} & 1 -\! 11\lambda +\! 3\lambda^2 +\! (1 -\! 2\lambda)\sqrt{-n} \\
24\lambda - 5 & 1 + \sqrt{-n} & 2 - 6\lambda & 3\lambda - \lambda\sqrt{-n} & 2\lambda + 3\lambda^2 + \lambda\sqrt{-n} \\
24\lambda - 13 & 1 + \sqrt{-n} & 4 - 6\lambda & -1 + 3\lambda + (1 - \lambda)\sqrt{-n} & 4 - 8\lambda + 3\lambda^2 + \lambda\sqrt{-n} \\
24\lambda - 17 & 1 + \sqrt{-n} & 5 - 6\lambda + \sqrt{-n} & 7 - 9\lambda + (1 - \lambda)\sqrt{-n} & 4 -\! 7\lambda +\! 3\lambda^2 +\! (2 -\! 2\lambda)\sqrt{-n}
\end{array}
\]

Ainsi on voit que pour tout nombre premier $n = 24\lambda - 1$, les premiers
termes des valeurs de Y et de Z seront
\[
\begin{array}{l}
Y = 2x^m + x^{m-1} + (1 - 6\lambda)x^{m-2} + (1 - 9\lambda)x^{m-3} + (1 - 11\lambda + 3\lambda^2)x^{m-4} + \text{etc.} \\
Z = x^{m-1} + x^{m-2} + (1 - \lambda)x^{m-3} + (1 - 2\lambda)x^{m-4} + \text{etc.}
\end{array}
\]

\page{132}
\note{page imprimée, titre courant « SUR UNE QUESTION D'ANALYSE. » et
numéro de page 17 ; le nombre 64 à l'encre en haut à droite ; au pied de la
page, la signature d'imprimeur « 3 ».}

Et dans cette formule on remarquera que les quatre premiers coefficients de
Y, savoir 2, 1, $1 - 6\lambda$, $1 - 9\lambda$, sont toujours moindres que
$\frac{1}{2}n$, mais le coefficient suivant $1 - 11\lambda + 3\lambda^2$, qui
est de l'ordre $3\lambda^2$ ou $\frac{n^2}{192}$, devient bientôt plus grand
que $n$ ; ainsi faisant $\lambda = 13$ ou $n = 311$, on a
$1 - 11\lambda + 3\lambda^2 = 365$.

Ces deux premiers tableaux pour les nombres premiers $4i + 1$ et $4i - 1$,
ne s'étendent pas au-delà de $A_4$ qui donne les coefficients de $x^{m-4}$
dans Y et Z ; pour aller plus loin, il faut subdiviser chacune des valeurs
$n = 24\lambda + 1$, 5, 13, 17, en quatre autres par la substitution des
cinq valeurs $\lambda = 5\mu$, $5\mu + 1$, $5\mu + 2$, $5\mu + 3$, $5\mu + 4$,
dont une doit être rejetée, comme ne donnant pas pour $n$ des nombres
premiers ; voici ces subdivisions :
\[
\begin{array}{c|l}
\textit{Valeur principale de } n. & \textit{Ses quatre sous-divisions.} \\ \hline
24\lambda + 1 & 120\mu + 1,\ 120\mu + 49,\ 120\mu + 73,\ 120\mu + 97 \\
24\lambda + 5 & 120\mu + 29,\ 120\mu + 53,\ 120\mu + 77,\ 120\mu + 101 \\
24\lambda + 13 & 120\mu + 13,\ 120\mu + 37,\ 120\mu + 61,\ 120\mu + 109 \\
24\lambda + 17 & 120\mu + 17,\ 120\mu + 41,\ 120\mu + 89,\ 120\mu + 113
\end{array}
\]

Appliquant notre méthode à ces différents cas, on trouvera pour chacun d'eux
l'expression générale du coefficient $A_5$ et celle de $A_6$, l'une
renfermant des termes affectés de $\mu^2$ et $\mu^2\sqrt{n}$, l'autre des
termes affectés de $\mu^2\sqrt{n}$ et $\mu^3$ ; on en conclura que le
coefficient de $x^{m-6}$ dans Y, étant de l'ordre $\mu$ deviendra bientôt
plus grand que $n^2$, ce qui manifeste l'augmentation progressive des
coefficients de Y. On voit en même temps que le nombre et la complication
des formules augmen-
\note{« de l'ordre $\mu$ » : la lettre, en fin de ligne, ne porte aucun
exposant lisible sur la vue.}

\page{133}
\note{page imprimée, titre courant « MÉMOIRE » et numéro de page 18, collée
au verso du feuillet dont le recto porte la page 17 (vue 132). Le « 64 » à
rebours en haut à gauche est le nombre du recto vu en transparence.}

tent rapidement avec la valeur de $n$, ce qui les rend peu propres à la
solution des cas particuliers. C'est pourquoi il suffira de développer ici
les formules relatives à l'une de nos subdivisions ; nous choisirons pour
cet effet la valeur $n = 120\mu + 61$.

On trouve alors, comme dans le cas de $n = 24\lambda + 13$ dont cette
subdivision fait partie
\[
S_1 = S_3 = S_4 = -\frac{1}{2} - \frac{1}{2}\sqrt{n},\quad S_2 = -\frac{1}{2} + \frac{1}{2}\sqrt{n}.
\]
On aura de plus $\left(\frac{5}{n}\right) = \left(\frac{61}{5}\right) = 1$,
$\left(\frac{6}{n}\right) = 1$, ce qui donne $S_5 = S_1$ et $S_6 = S_2$.

Cela posé dans le cas principal $n = 24\lambda + 13$, on a trouvé les
valeurs de $A_1$, $A_2$, $A_3$, $A_4$, dans lesquelles il faudra substituer
la valeur $\lambda = 5\mu + 2$, ce qui donnera, pour la division
$n = 120\mu + 61$, les valeurs
\[
\begin{array}{l}
A_1 = 1 + \sqrt{n} \\
A_2 = 16 + 30\mu \\
A_3 = -7 - 15\mu + (3 + 5\mu)\sqrt{n} \\
A_4 = 32 + 100\mu + 75\mu^2 - (2 + 5\mu)\sqrt{n}.
\end{array}
\]

Au moyen de ces valeurs on trouvera par les équations (1) :
\[
\begin{array}{l}
A_5 = -20 - \frac{195}{2}\mu - \frac{225}{2}\mu^2 + \left(6 + \frac{29}{2}\mu + \frac{15}{2}\mu^2\right)\sqrt{n} \\
A_6 = 63 + 246\mu + 280\mu^2 + 75\mu^3 - (3 + 14\mu + 15\mu^2)\sqrt{n},
\end{array}
\]
Et l'on voit que le terme $75\mu^3$ qui fait partie de $A_6$ deviendra plus
grand que $n^2$ si on a $n > \frac{(120)^3}{75}$ ou $n > 23040$.

Pour aller plus loin, c'est-à-dire pour trouver les valeurs

\page{134}
\note{page imprimée, titre courant « SUR UNE QUESTION D'ANALYSE. » et
numéro de page 19 ; le nombre 65 à l'encre en haut à droite. Les traits
pâles au-dessus de ce nombre sont le « 66 » écrit de l'autre côté du
feuillet (vue 135), vu en transparence.}

de $A_7$ et $A_8$, il faudrait au lieu de $\mu$ substituer l'une des six
valeurs $7\nu$, $7\nu + 1$, $7\nu + 3$, $7\nu + 4$, $7\nu + 5$, $7\nu + 6$,
c'est-à-dire il faudrait subdiviser la forme $n = 120\mu + 61$ en six autres
\[
n = 840\nu + 61,\ 181,\ 421,\ 541,\ 661,\ 781 ;
\]
alors $A_7$ contiendrait un terme de l'ordre $\mu^3\sqrt{n}$, et $A_8$ un
terme de l'ordre $\mu^4$. Il est inutile d'entrer dans de plus grands
détails ; nous déduirons seulement des valeurs précédentes de
$A_1$, $A_2 \ldots A_6$ celles des fonctions Y et Z développées jusqu'à la
puissance $x^{m-6}$ ; elles donnent pour le nombre premier $n = 120\mu + 61$,
\[
\begin{array}{l}
Y = 2x^m + x^{m-1} + (16 + 30\mu)x^{m-2} - (7 + 15\mu)x^{m-3} \\
\qquad + (32 + 100\mu + 75\mu^2)x^{m-4} - \left(20 + \frac{195}{2}\mu + \frac{225}{2}\mu^2\right)x^{m-5} \\
\qquad + (63 + 246\mu + 280\mu^2 + 75\mu^3)x^{m-6} - \text{etc.} \\
Z = x^{m-1} + (3 + 5\mu)x^{m-3} - (2 + 5\mu)x^{m-4} \\
\qquad + \left(6 + \frac{29}{2}\mu + \frac{15}{2}\mu^2\right)x^{m-5} - (3 + 14\mu + 15\mu^2)x^{m-6} \\
\qquad - \text{etc.}
\end{array}
\]

On en déduirait pour le cas de $n = 120\mu - 61$, les formules suivantes
qui satisfont à l'équation $4X = Y^2 + nZ^2$,
\[
\begin{array}{l}
Y = 2x^m + x^{m-1} + (16 - 30\mu)x^{m-2} - (7 - 15\mu)x^{m-3} \\
\qquad + (32 - 100\mu + 75\mu^2)x^{m-4} - \left(20 - \frac{195}{2}\mu + \frac{225}{2}\mu^2\right)x^{m-5} \\
\qquad + (63 - 246\mu + 280\mu^2 - 75\mu^3)x^{m-6} - \text{etc.} \\
Z = x^{m-1} + (3 - 5\mu)x^{m-3} - (2 - 5\mu)x^{m-4} \\
\qquad + \left(6 - \frac{29}{2}\mu + \frac{15}{2}\mu^2\right)x^{m-5} - (3 - 14\mu + 15\mu^2)x^{m-6} - \text{etc.}
\end{array}
\]
Dans les deux cas on a fait $m = \frac{1}{2}(n - 1)$.
\note{un filet imprimé ferme la page et, avec elle, le mémoire : le reste
de la page est blanc.}

\page{136}
\note{feuillet de la main de l'auteur, à l'encre, d'une écriture courante mais
lisible, collé sur le feuillet de montage ; le nombre 67 à l'encre en haut à
droite, avec un crochet tracé au-dessous. Le cachet rond « Bibliothèque
impériale, MSS. » est apposé au milieu de la page, sous la figure.}

\emph{Idées à vérifier.}

On demande si dans la surface pesante et également épaisse l'équation
d'équilibre des différentes courbes produites par les sections normales
n'est pas celle d'autant de chaînettes différentes ?

On remarque d'abord, $R$ étant le rayon de la sphère de moyenne courbure que
l'équation d'équilibre de la surface est $ge + T.\frac{2}{R} = 0$ ; on demande
si la courbe produite par la section normale qui fait l'angle de $45^\circ$
avec les sections principales courbures changeroit de forme dans le cas où
retranchant le reste de la surface elle demeureroit isolée. Pour que cette
courbe \uncertain{conserve} en effet la forme qui lui appartient lorsqu'elle
fait partie de la surface il faut que le retranchement supposé diminue la
tension $T$ dans le rapport d'un à 2.
\note{« avec les sections principales courbures » : la ligne porte ces mots,
sans « de » ; le « s » de « sections » et celui de « courbures » sont des
queues de fin de mot, lues selon le sens. Le mot rendu par « conserve »
est tracé serré et ne se laisse pas lire lettre à lettre. Dans le titre,
« à vé » est souligné.}

Soit $A$——$C$——$B$, $ASB$ cette courbe dont les points d'appuis sont $S$ et
$A$, $B$, lorsqu'elle est isolée ; lorsqu'elle fait partie de la surface dont
le contour est appuyé en $B$ $B''$ $B''$, il est évident que la force de
l'appui $B$ est augmentée par celle des appuis $B'$ $B''$ $B'''$ \&c. Si on
suppose que l'appui $B$ est proportionnel à $BC$ et que l'appui $B'$ augmente
la force de l'appui $B$ d'une quantité proportionnelle à la petite ligne
$b'B'$. Il est clair qu'en menant les petites lignes $B''b''$, $B'''b'''$ \&c.
\uncertain{Cela} épuiseroit la ligne $CB$ et on concluroit que, la force du
premier appui étant proportionnelle à cette ligne, elle se trouveroit en
effet doublée. Mais cette manière d'envisager les choses est loin d'être
satisfesante en effet les forces $B'$ $B''$ $B'''$ \&c semblent devoir être
proportionnelles aux lignes $Cb'$ $Cb''$ $Cb'''$ \&c. qui sont les cosinus des
angles $CB'B$, $CB''B$, $CB'''B$, plutôt qu'à leur \uncertain{complémens}
$b'B$ \&c et je n'ai pas trouvé la manière de détruire cette objection

\emph{(Retournez)}
\note{figure à la plume, dans la ligne, entre « Soit » et « $ASB$ cette
courbe » : une droite horizontale $A$——$C$——$B$ ; au-dessous, une courbe
qui part de $A$, descend jusqu'à un point bas marqué $S$ et remonte à $B$ ;
de $C$ partent trois droites montant vers des points $B'$, $B''$, $B'''$
situés au-dessus de la droite, près de $B$, d'où tombent sur elle trois
perpendiculaires aux pieds marqués $b'$, $b''$, $b'''$. Les accents des
lettres sont petits et serrés sur la figure. À la ligne « en $B$ $B''$
$B''$ », les accents sont lus tels qu'ils paraissent, deux traits sur
chacune des deux dernières lettres ; la figure porte $B'$, $B''$, $B'''$.
« \uncertain{Cela} » est repassé d'un trait épais qui le rend douteux.
« (Retournez) », souligné, est au pied de la page : la suite est au verso,
vue 137.}

\page{137}
\note{verso du feuillet de la vue 136, de la même main et de la même encre ;
la page fait suite au « (Retournez) » du recto. Aucun nombre n'y est écrit.}

Mais voisi une autre manière d'envisager les choses

Supposons une surface également pesante dans tous ses points et appuyée
dans tous les points de son contour sur un chassis circulaire. La force
d'appui ou de tension $T$ qui soutient la surface pesante peut être
décomposée suivant deux plans normaux entr'eux et à la surface, et comme
dans les suppositions présentes tout est égal par rapport à ces deux plans,
les deux portions de forces qui agissent dans chacun d'eux sont aussi égales
entr'elles.

L'élément est soutenu par ces deux parties de forces et on a l'équation
$ge + T.\frac{2}{R}$, $g$ est la pesanteur soit $g'$ la pesanteur dans une
simple ligne courbe, la pesanteur aura une partie constante $g''$. Soit la
pesanteur dans la courbe $2g''$ dans une surface elle sera $4g''$ on aura
donc en faisant $T = 2T'$, $4g''e + 2T'.\frac{2}{R} = 0$ et par conséquent
\struck{dans} $g''e + T'\frac{1}{R} = 0$ pour l'équation de la courbe produite
par une section normale. Il résulteroit delà que cette courbe est une
chaînette.
\note{la lettre qui suit $g$ et $g''$ dans les équations est un $e$ nettement
formé, ici et à la vue 136 ; l'équation « $ge + T.\frac{2}{R}$ » est écrite
sans « $= 0$ ». Le « $4g''$ » de la fin de phrase est repassé à la plume, et
dans « $4g''e$ » les accents de $g$ sont serrés au point qu'on n'en
distingue pas sûrement deux.}

Mais si à présent \struck{\uncertain{les}} forces des appuis augmente dans
certains points du contour et diminue dans les autres nous pourrons admettre
que ces forces, qui luttent contre la pesanteur élèvent les points d'appuis
qui se trouvent compris dans une des sections normales ; il y en aura une
pour laquelle cette \uncertain{élévation} sera la plus grande ; dans d'autre
il y aura abbaissement et l'abbaissement le plus grand
\note{le mot biffé avant « forces » est couvert d'une croix. Le mot rendu par
« élévation » est tracé en abrégé et se lit plutôt « ébration » : la
lecture est donnée d'après le sens. La page s'arrête sur « le plus grand »,
sans ponctuation, au tiers inférieur du feuillet.}

\page{138}
\note{nouveau feuillet de la même main, collé sur le feuillet de montage ; le
nombre 68 en haut à droite, d'un trait fin et pâle. Le texte n'occupe que la
moitié supérieure de la page ; au-dessous, une tache ronde et rien d'autre.}

Dans la vue d'étendre les conséquences qu'on peut tirer de l'existence de la
sphère de moyenne courbure j'ai cherché de tout côté ce qui concerne la
chaînette à cause de son rapport avec la voûte qui se soutient d'elle même,
je n'ai trouvé que dans l'histoire des mathématiques l'encyclopédie et le
mémoire de M poisson sur la surface élastique.

Je trouve dans un article de d'alembert « si une courbe est pressée « en
chaque point, par une puissance qui soit perpendiculaire « à la courbe, pour
qu'il y ait équilibre, il faut que cette puissance « soit en raison inverse
des rayon de courbure de cette courbe ».

Je voyois que ce principe s'appliqueroit également à une surface en
substituant la somme des \struck{rayons} raisons inverses des rayons de
courbures principales en sorte qu'on pourroit énoncer ainsi la proposition
générale

« \emph{Si une surface est pressée en chaque point par} « \emph{une puissance
qui soit perpendiculaire à sa} \struck{\emph{surface}} courbure « \emph{pour
qu'il y ait équilibre, il faut que cette puissance} « \emph{soit en}
\struck{\emph{raison}} \emph{inverse double raison inverse du rayon} « \emph{de
courbure de la sphère de moyenne courbure de la surface.}

ou ce qui est plus clair

\emph{Si une surface est pressée en chaque point par une puissance qui soit
perpendiculaire à} l'élément de \emph{sa courbure, pour qu'il y ait équilibre
il faut que la moitié de cette puissance soit en raison inverse du rayon de sa
sphère de moyenne courbure}.
\note{les guillemets ouvrants sont répétés en tête de chaque ligne de la
citation de d'Alembert et de l'énoncé général, à la manière du temps ; ils
sont conservés à leur place. L'énoncé général est souligné ligne à ligne et
embrassé à gauche par un long trait vertical ; « courbure » est écrit
au-dessus de la fin de ligne, au-dessus de « surface » biffé. Dans « soit en
raison inverse double raison inverse », seul « raison » porte un trait qui le
traverse ; « inverse », qui le suit, n'est que souligné comme le reste de la
ligne, et n'est donc pas donné pour biffé. Le second énoncé est souligné
aussi ; « l'élément de » y est ajouté au-dessus de la ligne, entre « à » et
« sa », et un trait passe sur la fin de « élément », qui peut être un simple
trait de liaison.}

\page{140}
\note{feuillet de la même main, collé sur le feuillet de montage ; le nombre
69 en haut à droite, du même trait fin et pâle que le 68 de la vue 138. Le
texte s'arrête aux deux tiers de la page.}

Voila ce qui doit être examiné.

L'élément $dxdy\sqrt{1 + p^2 + q^2}$ de la surface est, \struck{\uncertain{je
crois}} celui du plan tangent, si le plan tangent se confond dans tous ses
points avec celui des $x$ $y$ l'élément se réduit à $dxdy$ :
$\frac{1}{\sqrt{1 + p^2 + q^2}}$ est le cosinus de l'angle que le plan
tangent fait avec celui des $x$, $y$ ; $dxdy$ est la projection sur le plan
des $x$, $y$ de l'élément $dxdy\sqrt{1 + p^2 + q^2}$. La projection de cet
élément est son produit par le cosinus de l'angle qu'il fait avec le plan de
$x$ $y$, il est donc
$dxdy.\sqrt{1 + p^2 + q^2}.\frac{1}{\sqrt{1 + p^2 + q^2}} = dxdy$. comme nous
l'avons dit

A présent voila ce qu'il faudrait savoir expliquer si on prend la variation
de $\sqrt{1 + p^2 + q^2}$ on a $\frac{p\delta p + q\delta q}{\sqrt{1 + p^2 +
q^2}}$ et en intégrant par partie $\frac{2}{R}\delta z$ ainsi la variation
\note{le mot biffé après « est, » est couvert d'un trait épais ; « je
crois » est proposé sous la rature avec doute. Dans « $\sqrt{1 + p^2 + q^2}$ »
les exposants sont parfois écrits comme une barre sur la lettre
(« $\bar p$ ») ; ils sont donnés comme carrés partout, ainsi que la ligne
les écrit en clair ailleurs. Le « $= dxdy$ » final est porté à droite, sur
la ligne de « l'avons dit », et repassé. La page s'arrête sur « ainsi la
variation », sans suite.}

\end{document}
